% GENERATED_BY: run_oc_core_monograph_translation_rework_dispatch_v1.py
% AI_ASSISTED_TRANSLATION_DRAFT: TRUE
% THEOREM_REVIEW_STATUS: APPROVED
% THEOREM_REVIEW_MODE: FORMAL_AI_GATE__CODEX_CLI_CHATGPT
% THEOREM_REVIEW_REVIEWER_ID: CODEX_CLI_CHATGPT__AUTO_DISPATCH_20260406
% THEOREM_REVIEW_AUTHORITY_CLASS: FINAL_ADVERSARIAL_EXTERNAL_LLM_REVIEWER
% THEOREM_REVIEWED_AT: 2026-04-14T13:30:48Z
% THEOREM_REVIEW_DOSSIER_REF: public evidence replay route
% SOURCE_LANGUAGE: EN
% TARGET_LANGUAGE: RU
% SOURCE_EN_REF: content/17_oc_core_1_3_reader_guide.tex
\section{Как читать Core 1.3.3}

Core 1.3.3 рассчитан сразу на несколько типов читателей, и структура тома прямо отражает этот факт, а не прогоняет всех читателей через одну риторическую воронку.
Монография должна быть понятной и редакторам с рецензентами, которым нужна чистая научная форма изложения, и формально ориентированным читателям, которым нужны точные определения и зависимости между теоремами, и широкой научной аудитории, которой перед абстракцией требуется терпеливое объяснение, и читателям из предметных областей, которым нужно понимать, что именно эта рамка доказывает и чего она о более поздних проекциях пока не доказывает.

\subsection{Уровни читательской аудитории}

\textbf{Редакторы и научные рецензенты.}
Самый быстрый маршрут с высоким уровнем доверия состоит в том, чтобы прочитать аннотацию, настоящее руководство для читателя, главы с основными результатами, главы с обсуждением и фальсифицируемостью, а затем главу аудита исходников 1.3 ближе к концу основного текста. Этот путь раскрывает научную заявку тома, структуру утверждений, содержащих теоремы, и причину, по которой выпуск 1.3 честнее, чем сжатый пакет, скрывающий неразрешённое наследование за стилем.

\textbf{Формально ориентированные читатели.}
Если главный интерес состоит в теоремной дисциплине, естественный путь проходит через аксиоматику, семантику операторов, корпус теорем, главы о коллапсе и возрождении и приложения. Этот маршрут удерживает определения, гипотезы и выводимые следствия видимыми в пределах одного и того же тома, так что читателю не приходится курсировать между небольшой внешней статьёй и длинным внутренним реестром лишь затем, чтобы проверить, использует ли теорема именно то, что, как в ней сказано, она использует.

\textbf{Широкая научная аудитория.}
Читатели, приходящие из теории систем, науки о сложности или соседних фундаментальных направлений, могут сначала прочитать введение, предысторию, модель, результаты, обсуждение и заключение, а затем использовать разделы более поздних материалов 1.3, служащие руководством для читателя, чтобы понять, какие части унаследованного корпуса OC представлены как устойчивое ядро, какие части ограничены, а какие по-прежнему носят перспективный характер, а не полностью защищены.

\textbf{Читатели из предметных областей.}
Более поздние главы об уровнях, пространствах, циклах, экспериментах, фальсифицируемости, ветвящейся топологии, дисциплинах, модулях и контуре наблюдаемой вселенной остаются частью тома, потому что эта рамка не является лишь лозунгом о коллапсе. Это структурная программа, научный смысл которой зависит от того, как её формальное ядро проецируется в реальные классы моделей и в эмпирические или операционные линии проверки.

\subsection{Почему монография по-прежнему необходима}

Краткая журнальная статья по ядру, извлечённая из Core 1.3.3, преследует иную цель.
Она предназначена для того, чтобы представить одну ограниченную теоремную программу в форме, которую журнальный редактор или недоброжелательный рецензент может быстро проверить. Это ценно, но недостаточно для роли канонического основного научного документа этой рамки. Фундаментальная теория становится менее заслуживающей доверия, а не более, когда основной документ заменяется кратким пакетом, опускающим дидактическую лестницу, аксиоматическую оболочку, унаследованную архитектуру глав и границу между тем, что является структурным ядром, а что относится к более поздней проекции.

По этой причине монография остаётся первичным источником. Она даёт читателю весь научный объект: словарь, теорию, теоремы, доказательства, иллюстрации, ограничения и специфический для 1.3 аппарат повторной валидации, объясняющий, как текущий выпуск соотносится с унаследованным корпусом Core 1.2 и с более поздними работами OC/the methodological staging framework.

\subsection{Что нового в логике выпуска 1.3}

Выпуск 1.3 не сводится к формуле ``Core 1.2 плюс новый PDF''.
Он вводит явно сформулированную дисциплину публикации с приоритетом исходного текста:

\begin{itemize}
\item полный корпус LaTeX остаётся канонической средой авторской работы, а не рассматривается как сырой архив, скрытый за гораздо меньшей внешней статьёй;
\item утверждения, содержащие теоремы, разделяются по их судьбе, так что читатель может видеть, что сохраняется в исходной формулировке, что пересматривается и ограничивается, а что удерживается вне позитивного корпуса более узких производных артефактов;
\item научная идентичность, выстроенная вокруг ORCID, заменяет прежнюю публикационную оболочку, завязанную на пересылку через email;
\item монография и журнальная статья по ядру чётко разведены, так что и владелец, и читатель понимают, какой артефакт является основным научным томом, а какой представляет собой сжатое извлечение для журнального контура.
\end{itemize}

\subsection{Как организован том}

Первый крупный блок тома сохраняет унаследованный корпус OC: введение, предысторию, формальную модель, результаты, обсуждение, заключение, рисунки, границы, пороги, полную иерархию K-уровней, операторы, коллапс и возрождение, ветвящуюся топологию, предметные расширения, фальсифицируемость и модульную архитектуру. Эти главы не являются фоновым шумом; именно они образуют формальное и дидактическое тело, которое изначально придало Core 1.2 масштаб монографии.

Второй крупный блок представляет собой специфичный для 1.3 слой исправления и публикации.
Он объясняет дисциплину фиксации свидетельств, встроенную в исходный корпус, разбиение теоремных утверждений по их судьбе, хронологию того, что изменилось между унаследованным выпуском и нынешней публикационной линией с приоритетом исходного текста, а также отношение между полной монографией, извлечённой журнальной статьёй по ядру и более широким внешним контуром системы в the release engineering pipeline.

\subsection{Замечание о правдивости}

Этот том намеренно устроен так, чтобы уменьшить два распространённых режима провала.
Первый состоит в раздутой риторике: называть нечто завершённым просто потому, что оно велико.
Второй состоит во вводящей в заблуждение аскетичности: называть нечто строгим просто потому, что оно мало.
Core 1.3.3 пытается избежать и того и другого.
Он остаётся достаточно большим, чтобы правильно обучать теории и нести её формальную структуру в одном месте, но при этом достаточно явно формулирует свои граничные условия, чтобы более узкие производные артефакты можно было оценивать без скрытых заимствований.

\clearpage
